\section{メートル尺度アライメントとレンダリング教師信号}
\label{sec:render_supervision}

データ生成器は本研究の主要な方法論的貢献である。その目的は任意の合成コートを
作ることではない。再構成された実在のテニス施設を、画像とラベルが監査可能な単一の
幾何学的正解情報源を共有する、座標系を認識したレンダラへと変換する。

\subsection{座標系と公開再構成境界}
\label{sec:frames}

\(\T{b}{a}\) を、座標系 \(a\) で表された点を座標系 \(b\) に写す同次変換と記す。
公開再構成エクスポートは、右手系の正規化シーン座標系
\(\Scene_n\) を用いる。カメラは、右向きを \(+x\)、下向きを \(+y\)、前向きを
\(+z\) とするOpenCV座標軸を用い、\(\T{\Scene_n}{\Cam}\) を公開する。
\(\Scene_n\) とメートル値を持つシーン座標系 \(\Scene_m\) の間には、互いに逆数となる
一様スケールをメートル尺度アダプタが定義する。各規定コート \(i\) はメートル単位の
右手系座標系 \(\Court_i\) を持ち、\(+X\) は右サイドライン方向、\(+Y\) は奥側の
ベースライン方向、\(+Z\) は上向きである。これらの規約を
\cref{tab:coordinate_contracts}にまとめる。

再構成プログラムはstructure from motion、Gaussianの学習、およびラスタライズを
担う。われわれは2つの公開コマンドのみを介してこのプログラムと通信する。一方は
検証済みの標準シーンをエクスポートし、もう一方は与えられたカメラをレンダリングする。
エクスポート検証器は、アライメントを開始する前に、スキーマ、カメラ識別子、shape、
dtype、有限値、内部パラメータ、行列式が $+1$ の回転行列、および相互に逆な変換を検査する。この境界に
より、データセットが明示されていないcheckpointや座標規約に依存することを防ぐ。

\begin{table}[t]
  \centering
  \caption{Coordinate frames.  All transforms are stored as proper homogeneous
  matrices; no quaternion convention is required by the data generator.}
  \label{tab:coordinate_contracts}
  \small
  \begin{tabularx}{\linewidth}{@{}p{0.15\linewidth}p{0.20\linewidth}X@{}}
    \toprule
    Frame & Units / axes & Contract \\
    \midrule
    NHT scene \(\Scene_n\) & Normalized, right-handed & Public reconstruction
      and render coordinates; Gaussian means and captured cameras live here. \\
    Metric scene \(\Scene_m\) & Metres, right-handed & Reciprocal uniform
      scaling of \(\Scene_n\); common carrier for all accepted courts. \\
    Court \(\Court_i\) & Metres; \(+X\) right sideline, \(+Y\) far baseline,
      \(+Z\) up & Regulation geometry and physical point indices. \\
    Camera \(\Cam\) & OpenCV; \(+x\) right, \(+y\) down, \(+z\) forward &
      \(\T{\Scene}{\Cam}\) is camera-to-scene; projection uses its inverse. \\
    Canonical court \(\Canon\) & Metres, target-relative & Court frame after
      camera-end choice \(S\in\{I,R_z(\pi)\}\); pose10D authority. \\
    Model image & Augmented pixels & KP logits, principal point, focal length,
      and reprojection residual are all expressed in this same frame. \\
    \bottomrule
  \end{tabularx}
\end{table}



\subsection{画像上のラインからメートル尺度のコート配置へ}
\label{sec:alignment}

\begin{figure}[t]
  \centering
  \resizebox{\linewidth}{!}{% メートル系アライメントの流れ。すべて本稿のために作成した TikZ 図である。
\begin{tikzpicture}[
  font=\sffamily\scriptsize,
  frame/.style={draw=black!28, rounded corners=2.5pt, fill=none,
    inner sep=3mm, line width=0.45pt},
  title/.style={font=\sffamily\bfseries\footnotesize, text=black!85},
  tinylabel/.style={font=\sffamily\fontsize{6.3}{7.2}\selectfont,
    text=black!70, align=center},
  fitcam/.style={regular polygon, regular polygon sides=3, minimum size=3.2mm,
    inner sep=0pt, draw=gsblue, fill=gsblue!20, line width=0.5pt},
  holdcam/.style={regular polygon, regular polygon sides=3, minimum size=3.2mm,
    inner sep=0pt, draw=poseorange, fill=poseorange!22, line width=0.5pt},
  arr/.style={-{Stealth[length=1.7mm,width=1.2mm]}, line width=0.65pt},
  dot/.style={circle, inner sep=0pt, minimum size=1.2mm},
  gate/.style={rounded corners=1.5pt, inner sep=2pt, draw=lossred,
    fill=lossred!8, font=\sffamily\fontsize{6.2}{7.2}\selectfont, align=center},
]

  % ---------- (a) カメラ証拠を fit / holdout に先に分ける ----------
  \begin{scope}[local bounding box=pa]
    \node[fitcam, shape border rotate=-90] (f1) at (0.05,2.05) {};
    \node[fitcam, shape border rotate=-90] (f2) at (1.05,2.25) {};
    \node[fitcam, shape border rotate=-90] (f3) at (2.05,2.05) {};
    \node[holdcam, shape border rotate=-90] (h1) at (1.05,0.05) {};
    \draw[fill=softgray!70, draw=black!30, line width=0.4pt]
      (-0.35,0.60) rectangle (2.45,1.55);
    \node[tinylabel] at (1.05,1.08) {コート面の画像領域};
    \draw[arr, draw=gsblue] (f1) -- (0.40,1.55);
    \draw[arr, draw=gsblue] (f2) -- (1.05,1.55);
    \draw[arr, draw=gsblue] (f3) -- (1.70,1.55);
    \draw[arr, draw=poseorange, dashed] (h1) -- (1.05,0.60);
    \node[tinylabel, text=gsblue, anchor=south] at (1.05,2.50)
      {fit カメラ（変換の推定）};
    \node[tinylabel, text=poseorange, anchor=north] at (1.05,-0.25)
      {holdout カメラ（凍結検証）};
  \end{scope}

  % ---------- (b) 地面推定と光線交差 ----------
  \begin{scope}[shift={(4.6,0)}, local bounding box=pb]
    \coordinate (bl) at (0,-0.30);
    \coordinate (br) at (3.3,-0.30);
    \coordinate (tl) at (0.55,1.05);
    \coordinate (tr) at (2.75,1.05);
    \fill[softgray!80] (bl) -- (br) -- (tr) -- (tl) -- cycle;
    \draw[black!35, line width=0.4pt] (bl) -- (br) -- (tr) -- (tl) -- cycle;
    \draw[black!18] (0.50,-0.30) -- (0.95,1.05)
      (1.15,-0.30) -- (1.45,1.05)
      (1.85,-0.30) -- (1.90,1.05)
      (2.55,-0.30) -- (2.35,1.05);
    \node[fitcam, shape border rotate=-90] (g1) at (0.15,2.20) {};
    \coordinate (e1) at (1.05,0.30);
    \coordinate (e2) at (1.95,-0.20);
    \draw[arr, draw=gsblue] (g1) -- (e1);
    \draw[arr, draw=gsblue, dashed] (g1) -- (e2);
    \node[dot, fill=gsblue] at (e1) {};
    \node[dot, fill=gsblue] at (e2) {};
    \node[tinylabel, text=gsblue, anchor=south west] at (0.45,2.35)
      {線画素の光線と\\推定地上面の交差};
    \node[tinylabel, anchor=north, text width=42mm] at (1.65,-0.55)
      {採用画素数と光線--平面角が\\閾値未満なら、そのカメラを外す};
  \end{scope}

  % ---------- (c) 共通尺度の 2 面当てはめ ----------
  \begin{scope}[shift={(9.7,0)}, local bounding box=pc]
    \draw[courtgreen, line width=0.6pt, fill=courtgreen!7]
      (0.20,0.85) rectangle (1.75,2.05);
    \draw[courtgreen, line width=0.35pt] (0.975,0.85) -- (0.975,2.05);
    \draw[black!28, line width=0.35pt] (0.20,1.45) -- (1.75,1.45);
    \draw[courtgreen, line width=0.6pt, fill=courtgreen!7]
      (2.50,0.55) rectangle (4.05,1.75);
    \draw[courtgreen, line width=0.35pt] (3.275,0.55) -- (3.275,1.75);
    \draw[black!28, line width=0.35pt] (2.50,1.15) -- (4.05,1.15);
    \draw[<->, courtgreen, line width=0.5pt] (0.35,2.20) -- (3.90,2.20);
    \node[tinylabel, text=courtgreen, anchor=south] at (2.125,2.22)
      {共通尺度 \(s\)（メートル換算）};
    \node[gate, anchor=north] at (2.125,0.40)
      {中心間距離 \(\geq 10.97\) m、\\フットプリント重なり \(\leq 10^{-9}\)};
    \node[tinylabel, anchor=north, text width=44mm] at (2.125,-0.55)
      {トポロジーと尺度が食い違えば\\シーン全体を棄却する};
  \end{scope}

  % ---------- パネル枠と見出し ----------
  \node[frame, fit=(pa)] (fa) {};
  \node[frame, fit=(pb)] (fb) {};
  \node[frame, fit=(pc)] (fc) {};
  \node[title, anchor=south] at ($(fa.north)+(0,1.3mm)$) {(a) 証拠の分割};
  \node[title, anchor=south] at ($(fb.north)+(0,1.3mm)$) {(b) 地上面への持ち上げ};
  \node[title, anchor=south] at ($(fc.north)+(0,1.3mm)$) {(c) 共通尺度の 2 面当てはめ};

  \draw[arr] (fa.east) -- (fb.west);
  \draw[arr] (fb.east) -- (fc.west);
\end{tikzpicture}
}
  \caption{\textbf{フェイルクローズな3DGSからコートへのアライメント。}
  撮影カメラはフィッティング前に分割する。フィッティング用カメラのコートライン画素と
  ロバストな地面平面との交点を求める。2つの規定テンプレートを、シーン単位/メートルの
  単一の共有スケールでフィッティングした後、メートル尺度のシーン座標で剛体
  フィッティングを行う。固定した変換は、ホールドアウトカメラとコート全体のトポロジーを
  説明できなければならない。カメラ側エンドの正準化では、反射を含まない二つの回転 \(I\) と
  \(R_z(\pi)\) のみを用いる。}
  \label{fig:alignment}
\end{figure}

\paragraph{地面平面とリフトされた観測情報。}
疎なシーン点から地面支持点の候補を得る。RANSACにより
\(
  \pi_g=\{\mathbf{x}:\mathbf{n}^{\top}\mathbf{x}+b=0\}
\)
を推定し、その支持集合を3回精緻化したうえで、法線、支持集合、境界、またはカメラ側の
検査が妥当でない平面を棄却する。検出されたライン画素
\(\tilde{\mathbf{u}}\) に対し、その撮影カメラはシーン内のレイ
\(
  \mathbf{r}(\lambda)=\mathbf{o}+\lambda\mathbf{d}
\)
を定める。交点
\begin{equation}
  \lambda^*=-\frac{\mathbf{n}^{\top}\mathbf{o}+b}
                    {\mathbf{n}^{\top}\mathbf{d}}
  \label{eq:ray_plane}
\end{equation}
は、レイと平面の角度、正の距離、および施設境界に関する厳格なゲートをすべて通過した
場合にのみ保持する。観測可能なライン情報を持たないカメラは明示的に除外する。残る
観測情報が不十分な場合はフォールバックせず、エラーとする。

\paragraph{共有スケールを持つテンプレート仮説。}
リフトされた点群は、地面平面上の2次元基底で計測された観測集合
\(\mathcal{E}\) を形成する。メートル値を持つ正準テンプレート
\(\mathcal{Q}\) の各規定ライン区間をサンプリングする。
\(d(\mathbf{a},\mathcal{E})=\min_{\mathbf{e}\in\mathcal{E}}
\|\mathbf{a}-\mathbf{e}\|_2\) とすると、実装された提案スコアは
\begin{equation}
  J(\mathbf{c},\theta,s)=\frac{1}{|\mathcal{Q}|}
  \sum_{\mathbf{q}\in\mathcal{Q}}
  \exp\!\left[-\frac{1}{2}
    \left(\frac{d(s\mathbf{R}(\theta)\mathbf{q}+\mathbf{c},\mathcal{E})}
                    {s\sigma}\right)^2\right],
  \qquad
  (\mathbf{c}^*,\theta^*,s^*)=\arg\max J,
  \label{eq:template_fit}
\end{equation}
である。ここで、\(s\) はメートル当たりのNHTシーン単位で測られ、設定されたスコア幅は
\(\sigma=0.5\) mである。範囲を制限したdifferential evolution探索により、タイル状に
配置した中心、方位帯、およびスケールを網羅する。第1コートをフィッティングした後、
その観測情報を抑制してから第2仮説を探索する。2つの仮説は、ラインのトポロジーに
整合性があり、フットプリントの重複が無視でき、中心間距離が十分であり、かつ各々の
ネイティブスケールが宣言された許容範囲内で一致しなければならない。その後、単一の
共通スケールのもとで両者の中心を再フィッティングする。この複数コート制約により、
剛体のコート変換をフィッティングする前にメートル尺度を決定する。

受理された共通スケールを \(s_*\) とする。メートル尺度アダプタは
\begin{equation}
  \mathbf{x}_{\Scene_n}=s_*\mathbf{x}_{\Scene_m},\qquad
  \mathbf{x}_{\Scene_m}=s_*^{-1}\mathbf{x}_{\Scene_n}.
  \label{eq:metric_adapter}
\end{equation}
である。対応付けられた正準コート点 \(\mathbf{x}_j\) と、リフトされたメートル尺度の
シーン点 \(\mathbf{y}_j\) に対し、最終的な変換はKabsch解
\begin{equation}
  (\mathbf{R}_i,\mathbf{t}_i)=
  \arg\min_{\mathbf{R}\in\SOthree,\mathbf{t}\in\R^3}
    \sum_j \left\|\mathbf{R}\mathbf{x}_j+\mathbf{t}-\mathbf{y}_j\right\|_2^2,
  \quad
  \T{\Scene_m}{\Court_i}=
  \begin{bmatrix}\mathbf{R}_i&\mathbf{t}_i\\\mathbf{0}^{\top}&1\end{bmatrix}.
  \label{eq:kabsch}
\end{equation}
である。\(\det\mathbf{R}_i=1\) となるよう反射を補正する。スケールは
\cref{eq:kabsch}には現れない。これはすでに\cref{eq:metric_adapter}で決定済みであるためである。

\paragraph{分離された検証。}
カメラの所属はフィッティング前に固定する。フィッティング用カメラからの対応点で
\cref{eq:kabsch}を推定し、ホールドアウトカメラでは再フィッティングせずにその固定変換を
評価する。さらに、最近傍の対応点だけでなくコートテンプレート全体をサンプリングすることで、
可視な1本のベースライン上の残差が良好であっても、コート全体が誤っている状態を見逃さない。
区間ごとのcoverage、2コートのトポロジー、スケールの一致、および相互逆変換の検査が
すべて合格しなければならない。\cref{tab:alignment_gates}で選択したproduction閾値は、
報告されたベンチマーク精度ではなく、実行可能な受理条件である。完全なtyped configurationが
引き続き正とする情報源である。

\begin{table}[t]
  \centering
  \caption{\textbf{採用したアライメント判定の閾値。}
  これらは実行可能な受理条件であり、
  達成された精度ではない。実際の権威は型付き設定ファイル側にあり、
  本表はその要約である。
  完全な一覧は付録~\ref{app:thresholds} に置く。
  }
  \label{tab:alignment_gates}
  \small
  \begin{tabularx}{\linewidth}{@{}p{0.32\linewidth}X@{}}
    \toprule
    判定 & 閾値 \\
    \midrule
    証拠の所有権 & 決定的に選ぶ 48 カメラ。
      fit 8 : holdout 4 の単位で固定分割（48 カメラでは 32 / 16） \\
    地上面 & 支持点 \(\geq 1000\)、
      normal と上方向の余弦 \(\geq 0.97\)、
      正のカメラ側割合 \(=1\) \\
    投影証拠 & 有効なカメラあたり \(\geq 20\) 点 \\
    コート仮説 & ちょうど 2 面。中心間距離 \(\geq 10.97\) m、
      フットプリント重なり \(\leq 10^{-9}\) \\
    共通尺度 & 元尺度の相対不一致 \(\leq 0.07290401\) \\
    対応構築 & 距離 \(\leq 0.25\) m、カメラあたり 3--200 対応 \\
    fit 受理 & カメラ \(\geq 6\)、対応 \(\geq 100\)、
      0.30 m 以内の内点率 \(\geq 90\%\)、RMS と q95 \(\leq 0.30\) m \\
    holdout 受理 & カメラ \(\geq 3\)、対応 \(\geq 80\)、
      0.30 m 以内の内点率 \(\geq 90\%\)、RMS と q95 \(\leq 0.30\) m \\
    全テンプレート & 0.30 m 以内の内点率 \(\geq 60\%\)、
      q95 \(\leq 1.50\) m、全意味セグメントの内点率 \(\geq 50\%\) \\
    変換と投影 & 逆変換の絶対許容 \(10^{-6}\)、
      投影の絶対許容 \(10^{-5}\) px \\
    \bottomrule
  \end{tabularx}
\end{table}


\subsection{制御された3次元カメラ軌道}
\label{sec:trajectories}

\begin{figure}[t]
  \centering
  \resizebox{\linewidth}{!}{% Controlled camera synthesis from path families to leakage-free partitions.
\begin{tikzpicture}[
  font=\sffamily\scriptsize,
  panel/.style={draw=black!28, rounded corners=2.5pt, fill=white,
    minimum height=50mm, inner sep=0pt, line width=0.45pt},
  title/.style={font=\sffamily\bfseries\footnotesize, text=black!84},
  tinylabel/.style={font=\sffamily\fontsize{6.35}{7.15}\selectfont,
    text=black!70, align=center},
  arr/.style={-{Stealth[length=1.8mm,width=1.25mm]}, line width=0.65pt},
  camera/.style={regular polygon, regular polygon sides=3, minimum size=3.4mm,
    inner sep=0pt, draw=gsblue, fill=gsblue!17, line width=0.5pt},
  sample/.style={circle, fill=gsblue, draw=white, line width=0.3pt,
    minimum size=1.9mm, inner sep=0pt},
  splitbox/.style={rounded corners=1.5pt, minimum width=31mm,
    minimum height=6.6mm, align=left, inner xsep=3pt, line width=0.5pt},
]
  % Four compact panels correspond to generation decisions made once per group.
  \node[panel, minimum width=43mm] (pa) at (2.15,0) {};
  \node[panel, minimum width=40mm] (pb) at (6.43,0) {};
  \node[panel, minimum width=42mm] (pc) at (10.65,0) {};
  \node[panel, minimum width=39mm] (pd) at (14.82,0) {};
  \node[title, anchor=north] at ($(pa.north)+(0,-2.2mm)$) {(a) Path family};
  \node[title, anchor=north] at ($(pb.north)+(0,-2.2mm)$) {(b) 3-D arc length};
  \node[title, anchor=north] at ($(pc.north)+(0,-2.2mm)$) {(c) Target + look-at};
  \node[title, anchor=north] at ($(pd.north)+(0,-2.2mm)$) {(d) Group-disjoint split};

  % (a) Horizontal footprints and vertical profiles are independently controlled.
  \node[tinylabel, anchor=west] at (0.35,1.52) {ground footprint};
  \draw[courtgreen, line width=0.75pt] (1.12,0.78) circle[radius=0.54];
  \fill[courtgreen] (1.12,0.78) circle[radius=0.75pt];
  \draw[gsblue, line width=0.75pt, rotate around={22:(3.08,0.78)}]
    (3.08,0.78) ellipse[x radius=0.78,y radius=0.42];
  \fill[gsblue] (3.08,0.78) circle[radius=0.75pt];
  \node[tinylabel, text=courtgreen] at (1.12,0.04) {circle};
  \node[tinylabel, text=gsblue] at (3.08,0.04) {ellipse};
  \draw[black!22] (0.36,-1.36) -- (4.00,-1.36);
  \node[tinylabel, anchor=west] at (0.72,-0.30) {camera height};
  \draw[courtgreen, line width=0.75pt] (0.58,-0.88) -- (1.88,-0.88);
  \draw[gsblue, line width=0.75pt]
    plot[smooth] coordinates {(2.22,-0.88) (2.48,-0.57) (2.75,-0.46)
      (3.03,-0.68) (3.30,-1.00) (3.58,-1.10) (3.88,-0.88)};
  \draw[arr, draw=black!48] (0.46,-1.31) -- (0.46,-0.55)
    node[above, tinylabel] {$z$};
  \node[tinylabel, text=courtgreen] at (1.23,-1.16) {planar};
  \node[tinylabel, text=gsblue] at (3.05,-1.39) {sinusoidal};
  \node[tinylabel, text width=37mm, anchor=south] at (2.15,-2.30)
    {$p(\theta)=[a\cos\theta,\,b\sin\theta,\,
      h+A\sin(k\theta+\phi)]^\top$};

  % (b) Samples are equally spaced along the full spatial curve, not in theta.
  \coordinate (o3) at (4.86,-1.18);
  \draw[arr, draw=black!45] (o3) -- ++(0.52,0) node[right, tinylabel] {$x$};
  \draw[arr, draw=black!45] (o3) -- ++(-0.27,0.34) node[above, tinylabel] {$y$};
  \draw[arr, draw=black!45] (o3) -- ++(0,0.60) node[above, tinylabel] {$z$};
  \draw[gsblue, line width=0.9pt]
    plot[smooth cycle, tension=0.55] coordinates {
      (5.10,-0.46) (5.31,0.16) (5.91,0.72) (6.73,0.98)
      (7.47,0.72) (7.86,0.14) (7.69,-0.49) (7.05,-0.89)
      (6.28,-0.96) (5.55,-0.80)};
  \foreach \x/\y in {5.10/-.46,5.31/.16,5.91/.72,6.73/.98,
      7.47/.72,7.86/.14,7.69/-.49,7.05/-.89,6.28/-.96,5.55/-.80}
    \node[sample] at (\x,\y) {};
  \draw[black!30, dashed] (5.91,0.72) -- (5.91,-1.22)
    (7.47,0.72) -- (7.47,-1.22);
  \draw[<->, >=Stealth, black!55, line width=0.5pt]
    (5.95,0.86) -- node[above, tinylabel] {$\Delta s$} (6.67,1.08);
  \node[tinylabel, fill=white, inner sep=1pt] at (6.43,-1.28)
    {equal spatial steps};
  \node[tinylabel, text width=35mm, anchor=south] at (6.43,-2.30)
    {$s(\theta)=\int_0^\theta\!\|p'(u)\|_2\,du$,\quad
     $s_i=iL/N$\\dense curve $\rightarrow$ invert cumulative $s$};

  % (c) A complex-centred sample binds to the metrically nearest court.
  \draw[courtgreen, fill=courtgreen!6, line width=0.48pt]
    (8.85,-0.82) rectangle (10.13,0.12);
  \draw[courtgreen, line width=0.42pt] (9.49,-0.82) -- (9.49,0.12);
  \draw[courtgreen, fill=courtgreen!13, line width=0.65pt]
    (11.18,-0.82) rectangle (12.46,0.12);
  \draw[courtgreen, line width=0.42pt] (11.82,-0.82) -- (11.82,0.12);
  \fill[courtgreen] (9.49,-0.35) circle[radius=0.8pt]
    (11.82,-0.35) circle[radius=0.8pt];
  \node[tinylabel] at (9.49,-1.06) {$\Court_A$};
  \node[tinylabel, text=courtgreen] at (11.82,-1.06) {$\Court_B=c_i^*$};
  \node[camera, shape border rotate=-75] (vc) at (11.35,1.13) {};
  \node[camera, shape border rotate=-42, opacity=0.55] (vp) at (10.60,0.86) {};
  \node[camera, shape border rotate=-122, opacity=0.55] (vn) at (12.24,0.78) {};
  \draw[gsblue!45, line width=0.55pt] (vp) to[bend left=10] (vc)
    to[bend left=10] (vn);
  \draw[black!42, dashed] (vc) -- node[left, tinylabel] {$d_A$} (9.49,-0.35);
  \draw[arr, draw=poseorange] (vc) -- node[right, tinylabel, text=poseorange]
    {$d_B$} (11.82,-0.35);
  \draw[arr, draw=gsblue!72] (vp) -- (11.82,-0.35);
  \draw[arr, draw=gsblue!72] (vn) -- (11.82,-0.35);
  \node[tinylabel, text width=37mm, anchor=south] at (10.65,-2.30)
    {$c_i^*=\operatorname*{lexmin}\arg\min_c
      \|v_i-o_c\|_2$\\camera $+z$ axes look at $o_{c_i^*}$};

  % (d) Seeded allocation moves complete trajectory groups as indivisible units.
  \node[tinylabel, draw=black!28, fill=softgray!72, rounded corners=1.5pt,
    minimum width=31mm, minimum height=8.0mm] (groups) at (14.82,1.08)
    {trajectory groups $G_1,\ldots,G_n$\\[-0.1em]seeded shuffle};
  \coordinate (forktop) at (13.10,0.70);
  \node[splitbox, draw=courtgreen, fill=courtgreen!9] (train) at (14.82,0.18)
    {\textcolor{courtgreen}{\textbf{Train 80\%}}\quad $G_2,G_4,\ldots$};
  \node[splitbox, draw=gsblue, fill=gsblue!8] (val) at (14.82,-0.58)
    {\textcolor{gsblue}{\textbf{Val. 10\%}}\quad $G_7$};
  \node[splitbox, draw=poseorange, fill=poseorange!8] (test) at (14.82,-1.34)
    {\textcolor{poseorange}{\textbf{Test 10\%}}\quad $G_9$};
  \draw[black!52, line width=0.55pt] (groups.south) -- ++(0,-0.12) --
    (forktop) -- (13.10,-1.34);
  \draw[arr, draw=courtgreen] (13.10,0.18) -- (train.west);
  \draw[arr, draw=gsblue] (13.10,-0.58) -- (val.west);
  \draw[arr, draw=poseorange] (13.10,-1.34) -- (test.west);
  \node[tinylabel, text width=33mm, anchor=south] at (14.82,-2.30)
    {all frames and views of one $G_i$ stay together};
\end{tikzpicture}
}
  \caption{\textbf{決定論的なビュー生成。}  円および楕円の候補軌道を撮影カメラのoffsetから
  スケーリングし、高さ方向には平面または正弦波を用いる。密な補間により3次元弧長に
  関して一様にサンプリングする。各カメラは決定論的に解決された対象コートを注視し、
  個々のフレームではなく軌道グループ全体を80/10/10のsplitに割り当てる。}
  \label{fig:trajectories}
\end{figure}

生成器はテニス施設に対して1つ、また受理された各コートに対して1つの周回中心を導出する。
基準半径には、その中心に対する撮影カメラの水平offsetの90パーセンタイルを用い、提案の
スケールを再構成された会場に固有のものとする。次に、型付き直積により、円/楕円の形状、
半径スケール、軸比、平面内方位、基準高度、および鉛直profileを列挙する。局所軌道は
\begin{align}
  \bm{\gamma}(\theta)_{xy}
    &=\mathbf{R}(\phi)
      \begin{bmatrix}a\cos\theta\\b\sin\theta\end{bmatrix},
      &a&=r_{90}\alpha, &b&=a\beta, \\
  \gamma_z(\theta)
    &=h_0+A\sin(n\theta+\varphi),
  \label{eq:trajectory}
\end{align}
の形をとる。\(\beta=1\) なら円、\(A=0\) なら平面軌道となる。すべての高さは
正でなければならない。

\(\theta\) の一様な増分は、楕円や鉛直変調された軌道上の距離に関して一様ではない。
そこで、4096点をサンプリングして累積3次元弧長関数を構成し、弧長について一様に補間する。
さらに、閉軌道上の隣接stepが設定された最大値以下になるまでサンプル数を増加させる。
候補の順序と選択には、固定seedと安定なfieldの明示的な順序を用いる。coverage objectiveは、
提案budgetの範囲内で、多様なcoverage/semantic-support modeと軌道グループを保持する。

以下に示すsingletonおよびカメラごとのtargetの挙動は、
\texttt{dataset/court=v2}によって明示的に有効化される。後方互換なパイプラインselectorは
v1（7つの2点channelおよびgroup-level target assignment）のままである。この生成スキーマv2は、
Court Detection V3における利用側／姿勢の契約とは別物である。v2はシリアライズされるsingleton
labelを定義するのに対し、V3はload後に導出される正準target-court training authorityを定義する。
v2では、コート中心軌道はその中心コートを維持する。施設中心軌道では、メートル尺度の
3次元空間で各カメラ中心に最も近い受理済みコートを選ぶ。距離が \(10^{-9}\) m以内で
同値の場合は辞書順で決定する。カメラ中心を \(\mathbf{c}\)、選ばれたコート上の注視点を
\(\mathbf{q}\) とし、次式によりOpenCVカメラ基底を構成する。
\begin{equation}
  \mathbf{f}=\frac{\mathbf{q}-\mathbf{c}}{\|\mathbf{q}-\mathbf{c}\|_2},\quad
  \mathbf{r}=\mathbf{f}\times\mathbf{u}_{\mathrm{world}},\quad
  \mathbf{d}=\mathbf{f}\times\mathbf{r},
  \label{eq:lookat}
\end{equation}
これにより、各軸は右、下、前を向く。水平画角を \(\omega\)、出力幅を \(W\) とすると、
\begin{equation}
  f_x=f_y=\frac{W}{2\tan(\omega/2)},\qquad
  (c_x,c_y)=\left(\frac{W-1}{2},\frac{H-1}{2}\right).
  \label{eq:intrinsics}
\end{equation}
完全な軌道をサンプリングしてからsplitを割り当てる。同じ軌道グループ識別子を共有する全
フレームをtrain、validation、testのいずれか1つのsplitに保持することで、ほぼ重複する
隣接ビューがsplit間にleakすることを防ぐ。

\begin{table}[t]
  \centering
  \caption{\textbf{軌道とサンプリングの方針。}
  数値は設定上の下限または予算であり、
  実際に採用された視点分布の証明ではない。
  採用後の分布は manifest から集計して報告する。
  }
  \label{tab:data_policy}
  \small
  \begin{tabularx}{\linewidth}{@{}p{0.26\linewidth}X@{}}
    \toprule
    次元 & 設定値 \\
    \midrule
    平面経路 & 円・楕円・矩形・超楕円。軸比 \(1.0,0.8,0.65\)。
      方位 \(0^\circ,45^\circ,90^\circ\) \\
    半径と高さ & 復元オフセットの \(0.65\)--\(0.95\) 倍。
      基準高さ \(1.5,2.25,3.0\) m \\
    鉛直経路 & 平面または正弦。振幅 \(0.0,0.25,0.5\) m \\
    視野 & 注視点高さ \(0\) または \(1.5\) m。
      水平画角 \(45^\circ\)--\(110^\circ\)（会場別の設定） \\
    サンプリング & seed 695。三次元弧長一様。隣接間隔の上限 1.05 m \\
    被覆 & 提案予算 4800。軌道グループ下限 24、採用フレーム下限 2000 \\
    分割と描画 & 軌道グループ単位で 8:1:1。8 shard \\
    SfM 制約 & 凸包を 0.5 m 内側へオフセット。空間被覆セル 1.0 m \\
    \bottomrule
  \end{tabularx}
\end{table}


\subsection{レンダリングと厳密なラベル}
\label{sec:labels}

メートル尺度のカメラを正規化シーン単位へ変換し、公開3DGSレンダラに渡す。レンダラはRGB、
depth、alpha、および画素support情報を返す。コートパイプラインは、別途レンダリングしたコートを
画像へ貼り付けない。そのため、背景、表面appearance、およびocclusion構造は元の動画から
再構成された状態を維持する。

正準な物理点 \(\mathbf{X}_k\) に対し、生成器は
\begin{equation}
  \mathbf{X}^{\Cam}_k =
  \T{\Cam}{\Scene_m}\T{\Scene_m}{\Court_i}\mathbf{X}_k,
  \qquad
  \tilde{\mathbf{u}}_k=\Kmat\mathbf{X}^{\Cam}_k,
  \qquad
  \mathbf{u}_k=\left(\frac{\tilde u_k}{\tilde w_k},
                       \frac{\tilde v_k}{\tilde w_k}\right).
  \label{eq:label_projection}
\end{equation}
を計算する。各recordは、メートル尺度のシーン位置、カメラdepth、画素/正規化座標、
in-front状態、in-frame状態、およびrenderer-support flagを保持する。このflagは、半径1画素の
近傍に十分なレンダリングalphaと正のレンダリングdepthが含まれることを意味するが、
メートル尺度のpoint-depth occlusion testではない。consumer contractは14個の物理点を
カメラ相対のnear/far singleton channelへ並べ替えるが、画像座標からidentityを推測することは
ない。segmentationおよびline targetは同じ投影コートgeometryから実体化する。学習時のcrop、
affine/perspective、color、blur、geometric validity、およびsupervision-eligibilityの処理は、
1つの拡張後の画像座標系において、教師信号を持つすべての量を更新する。

公開sampleはpose10Dフィールドを重複して保持せず、\(\T{\Scene_m}{\Cam}\)、
\(\T{\Scene_m}{\Court_i}\)、および \(\Kmat\) を格納する。画像拡張後、model adapterは
カメラlabel
\begin{equation}
  \T{\Court_i}{\Cam}=
  \bigl(\T{\Scene_m}{\Court_i}\bigr)^{-1}\T{\Scene_m}{\Cam}.
  \label{eq:camera_to_court}
\end{equation}
を導出する。これは3つの並進自由度と3つの回転自由度を持つ。次にadapterは
\cref{eq:canonical_pose}のtarget側canonicalizationを適用し、pose10D教師を構成する。
拡張後の焦点距離は剛体変換とは別に保持する。これはnetworkが再投影のために焦点距離を
同時に予測するためである。これが本論文における「6DoF supervision」の意味であり、
別taskで用いられるyawのみのplayer-location labelと混同してはならない。

あるシーンについて \cref{eq:camera_to_court} が受理されれば、その後の別ビューのレンダリングには、
点のclick、line tracing、PnP solve、フレームごとのextrinsic calibrationのいずれも必要ない。
人手は適切な動画の取得とシーン単位のアライメントの監査に集中する。この償却効果によって、
\cref{sec:model,sec:consistency}の大域姿勢目的関数を大規模に適用することが現実的になる。
